% book/chapters/ch09-context.tex
\chapter{Context Management}
\label{ch:context}

\index{context}\index{CLAUDE.md}

Every Claude Code session starts with a fresh context window. Two mechanisms
let knowledge persist across sessions:

\begin{itemize}
  \item \textbf{CLAUDE.md files} — instructions you write that Claude reads on
        every session start.
  \item \textbf{Auto memory} — notes that Claude writes for itself based on
        your corrections, then re-reads in future sessions.
\end{itemize}

This chapter covers both, plus the rule files and import system that organize
larger projects.

\section{CLAUDE.md Hierarchy}
\label{sec:claude-md-scopes}

\index{CLAUDE.md!scopes}

CLAUDE.md files load in order, from broadest scope to most specific. Later
files appear later in context, so project-level instructions override
user-level ones on conflict.

\begin{table}[ht]
\centering
\caption{CLAUDE.md locations and scope}
\label{tab:claude-md-scopes}
\begin{tabularx}{\textwidth}{llX}
\toprule
\textbf{Scope} & \textbf{Location} & \textbf{Use case} \\
\midrule
Managed policy & \texttt{/etc/claude-code/CLAUDE.md} (Linux/WSL)\newline
                 \texttt{/Library/Application Support/ClaudeCode/CLAUDE.md} (macOS)\newline
                 \texttt{C:\textbackslash{}Program Files\textbackslash{}ClaudeCode\textbackslash{}CLAUDE.md} (Windows)
               & Organization-wide standards deployed by IT \\
User           & \texttt{\textasciitilde/.claude/CLAUDE.md}
               & Personal preferences across every project \\
Project        & \texttt{./CLAUDE.md} or \texttt{./.claude/CLAUDE.md}
               & Team-shared, committed to version control \\
Local          & \texttt{./CLAUDE.local.md}
               & Personal project notes; add to \texttt{.gitignore} \\
\bottomrule
\end{tabularx}
\end{table}

When Claude Code starts in a subdirectory, it walks up the directory tree and
loads every \texttt{CLAUDE.md} it finds along the way. \texttt{CLAUDE.md}
files in subdirectories \emph{below} the working directory load on demand —
only when Claude reads a file in that subdirectory.

\begin{notebox}
Run \texttt{/init} in a new project to generate a starting \texttt{CLAUDE.md}.
Claude analyzes your codebase and proposes build commands, test instructions,
and conventions. If a \texttt{CLAUDE.md} already exists, \texttt{/init}
suggests refinements rather than overwriting.
\end{notebox}

\subsection{Size Targets}

Aim for under \textbf{200 lines per CLAUDE.md file}. The file is loaded into
context at every session start, so every kilobyte costs tokens forever.
Longer files also reduce adherence — when the file gets too long, Claude
starts ignoring lines further from the top.

\begin{lstlisting}[caption={A focused CLAUDE.md (\textasciitilde 30 lines)}]
# Code style
- Use ES modules (import/export), not CommonJS (require)
- 2-space indentation, double quotes for strings
- Prefer functional components with hooks

# Workflow
- Typecheck after a series of edits: `npm run typecheck`
- Run a single failing test, not the whole suite, during iteration
- Open the PR with `gh pr create` once tests pass

# Don't
- Don't add comments that just restate the code
- Don't introduce new dependencies without asking
\end{lstlisting}

\subsection{Imports and Composition}

\index{CLAUDE.md!imports}

CLAUDE.md files can import other Markdown using \texttt{@path/to/file}
syntax. Imports expand at session start (up to a depth of four hops).

\begin{lstlisting}[caption={Importing related files into CLAUDE.md}]
See @README.md for the project overview and @package.json
for the available scripts.

# Additional instructions
- Git workflow: @docs/git-instructions.md
- Personal overrides: @~/.claude/my-project-instructions.md
\end{lstlisting}

For private, machine-local notes, create \texttt{CLAUDE.local.md} at the
project root and add it to \texttt{.gitignore}.

\begin{exercisebox}[Exercises — building your project's CLAUDE.md]
These exercises continue the running project introduced in earlier chapters.
You should have a \texttt{.claude/} directory at your project root.

\begin{enumerate}
  \item Run \texttt{/init} inside the project. Inspect the generated
        \texttt{CLAUDE.md}: how many lines does it produce, and which
        instructions did Claude infer from your codebase?
  \item Trim the generated file to under 100 lines. Remove anything that
        Claude could derive by reading the code (file-by-file
        descriptions, language conventions, etc.). Keep only build
        commands, project-specific rules, and ``don'ts.''
  \item Add a \texttt{CLAUDE.local.md} with one private preference (e.g.,
        a sandbox URL or a per-developer dev port). Confirm
        \texttt{git status} does \emph{not} list it as untracked — i.e.,
        it is already in your \texttt{.gitignore}.
  \item Start a new session and run \texttt{/memory}. Verify that both
        \texttt{CLAUDE.md} and \texttt{CLAUDE.local.md} are listed as
        loaded.
\end{enumerate}
\end{exercisebox}

\section{Path-Scoped Rules}
\label{sec:claude-rules}

\index{CLAUDE.md!rules}

Once a CLAUDE.md grows past 200 lines, split topical instructions out into
\texttt{.claude/rules/}. Rules can be scoped to specific file paths via YAML
frontmatter so they only load when Claude works on matching files:

\begin{lstlisting}[language=bash,caption={Layout of \texttt{.claude/rules/}}]
.claude/
  CLAUDE.md                # baseline project instructions
  rules/
    api-design.md          # always-loaded
    frontend/
      react-conventions.md # scoped to src/components/**
    backend/
      sql-style.md         # scoped to db/migrations/**
\end{lstlisting}

\begin{lstlisting}[caption={A path-scoped rule}]
---
paths:
  - "src/api/**/*.ts"
  - "src/api/**/*.tsx"
---

# API conventions

- Validate every request body with a Zod schema
- Wrap handlers in `asyncHandler` for error propagation
- Return `{ success, data?, error? }` from every endpoint
\end{lstlisting}

A rule without a \texttt{paths} field loads unconditionally alongside
\texttt{CLAUDE.md}.

\section{Auto Memory}
\label{sec:auto-memory}

\index{auto memory}

Auto memory lets Claude accumulate notes between sessions without you writing
anything. When you correct Claude, fix a misunderstanding, or teach it a
build command, Claude decides whether the lesson is worth keeping; if so, it
writes a Markdown file in:

\begin{lstlisting}[language=bash,caption={Auto-memory storage location}]
~/.claude/projects/<project>/memory/
  MEMORY.md          # concise index, loaded into every session
  debugging.md       # detailed notes Claude wrote
  api-conventions.md
  ...
\end{lstlisting}

Only the first 200 lines (or 25\,KB) of \texttt{MEMORY.md} load at session
start. Topic files load on demand. The directory is per git repository and
is shared across worktrees on the same machine.

\subsection{Inspecting and Editing Memory}

Run \texttt{/memory} in a session to:
\begin{itemize}
  \item see which CLAUDE.md and rule files are currently loaded,
  \item toggle auto memory on or off,
  \item open the memory directory in your editor.
\end{itemize}

Memory files are plain Markdown — edit or delete them with no special tooling.

\subsection{Disabling Auto Memory}

Auto memory is on by default. To turn it off:

\begin{lstlisting}[language=json,caption={Disabling auto memory in \texttt{.claude/settings.json}}]
{ "autoMemoryEnabled": false }
\end{lstlisting}

Or set \texttt{CLAUDE\_CODE\_DISABLE\_AUTO\_MEMORY=1} in your environment.

\begin{warningbox}
Auto memory is machine-local: it is \emph{not} synced across machines,
laptops, or cloud sandboxes. If you change machines, Claude starts the new
environment without the accumulated notes.
\end{warningbox}

\section{Compaction and What Survives}
\label{sec:compaction}

\index{compaction}

When the active context window approaches its limit, Claude Code automatically
\emph{compacts}: it summarizes older turns to free space while preserving
the key code, files, and decisions. You can also trigger it manually:

\begin{itemize}
  \item \texttt{/clear} — full reset; nothing from the previous turns survives.
  \item \texttt{/compact} — summarize the entire conversation.
  \item \texttt{/compact <instructions>} — summarize with a focus, e.g.\
        \texttt{/compact focus on the API changes}.
  \item \texttt{Esc}, \texttt{Esc} — open the rewind menu, pick a message,
        and either restore from that point or summarize around it.
\end{itemize}

\begin{notebox}
Project-root \texttt{CLAUDE.md} \emph{is} re-injected after compaction.
Nested \texttt{CLAUDE.md} files in subdirectories are \emph{not} — they
reload the next time Claude reads a file in that subdirectory. If an
instruction needs to persist across compactions, put it in the root
\texttt{CLAUDE.md}, not in a nested one.
\end{notebox}

\section{Monorepos and Excludes}
\label{sec:monorepos}

In a large monorepo, ancestor \texttt{CLAUDE.md} files from other teams may
be irrelevant or even contradict your team's conventions. Exclude them with
\texttt{claudeMdExcludes}:

\begin{lstlisting}[language=json,caption={Excluding ancestor CLAUDE.md files}]
{
  "claudeMdExcludes": [
    "**/monorepo/CLAUDE.md",
    "/home/me/monorepo/other-team/.claude/rules/**"
  ]
}
\end{lstlisting}

Add this to \texttt{.claude/settings.local.json} so the exclusion stays
on your machine. Managed-policy \texttt{CLAUDE.md} files cannot be excluded.

\begin{exercisebox}[Exercises — rules, memory, compaction]
These exercises build on the previous set.

\begin{enumerate}
  \item Take a 100-line section of your trimmed \texttt{CLAUDE.md} that
        applies to \emph{only} one subdirectory (e.g.\ \texttt{src/api/}
        or \texttt{tests/}). Move it into a path-scoped rule under
        \texttt{.claude/rules/}, with a \texttt{paths:} pattern matching
        that subdirectory. Confirm \texttt{/memory} only lists the rule
        when you reference a file inside that directory.
  \item Provoke an auto-memory write: in a session, correct Claude on a
        consistent project convention twice in a row (e.g.\
        ``always use pnpm, not npm''). End the session. Open the
        \texttt{\textasciitilde/.claude/projects/<project>/memory/}
        directory and inspect what was written. Edit or delete entries
        that aren't useful.
  \item In a long session, run \texttt{/compact focus on the auth
        changes}. Verify the next message still has access to the auth
        files; ask Claude to list which files it remembers editing.
  \item In a monorepo, add a \texttt{claudeMdExcludes} entry to
        \texttt{.claude/settings.local.json} that filters out one other
        team's directory. Confirm \texttt{/memory} no longer lists files
        from that path.
\end{enumerate}
\end{exercisebox}

\section{AGENTS.md Compatibility}
\label{sec:agents-md}

If your repository already has an \texttt{AGENTS.md} for another agent
(Cursor, Continue, etc.), Claude Code does \emph{not} read it directly — but
you can import it from your \texttt{CLAUDE.md}:

\begin{lstlisting}[caption={Sharing instructions between agents}]
@AGENTS.md

## Claude Code

Use plan mode for changes under `src/billing/`.
\end{lstlisting}

A symlink (\texttt{ln -s AGENTS.md CLAUDE.md}) also works on macOS and
Linux. On Windows, use the \texttt{@AGENTS.md} import — symlinks require
Administrator privileges or Developer Mode.

%% ---------------------------------------------------------------
\begin{summarybox}
\begin{itemize}
  \item Persistent instructions live in \texttt{CLAUDE.md} at four scopes:
        \textbf{managed}, \textbf{user}, \textbf{project}, and \textbf{local};
        every level loads, with later levels overriding earlier ones on
        conflict.
  \item Target \textbf{under 200 lines} per CLAUDE.md. Use
        \texttt{@path/to/file} imports for composition and
        \texttt{.claude/rules/} for path-scoped topical instructions.
  \item \textbf{Auto memory} is enabled by default. Claude writes per-project
        notes to
        \texttt{\textasciitilde/.claude/projects/<project>/memory/}; the first
        200 lines of \texttt{MEMORY.md} load every session.
  \item Project-root \texttt{CLAUDE.md} survives \texttt{/compact};
        nested ones reload on demand. \texttt{/clear} discards everything.
  \item In monorepos, use \texttt{claudeMdExcludes} to skip other teams'
        files. Managed-policy CLAUDE.md cannot be excluded.
\end{itemize}
\end{summarybox}

\printbibliography[heading=subbibliography,title={References}]
